\section{Cómo hacerlo crecer}\label{sec:crecer}

KalaScript no está terminado ni debe estarlo: tiene que crecer con mis apuntes
reales. Esta sección fija cómo se añade, cambia o retira un signo sin que los
apuntes viejos dejen de leerse.

\subsection{Leyenda personal}\label{sec:leyenda}
La primera página de cada cuaderno es la leyenda: las abreviaturas propias, los
nombres frecuentes del curso y los signos candidatos que estoy probando. Si
alguien (o yo dentro de un año) abre el cuaderno, ahí está todo lo que no
aparece en este manual.

\subsection{Cuándo merece un signo nuevo}
Durante la clase no invento nada: escribo la palabra y la subrayo con una línea
ondulada. Al repasar, cuento cuántas veces la subrayé. Una idea pasa a
candidata si cumple las cuatro condiciones:
\begin{enumerate}
\item \textbf{Frecuencia.} Aparece varias veces por semana, en más de un curso.
\item \textbf{Ganancia.} El signo es claramente más rápido o más preciso que
 la palabra o su abreviatura.
\item \textbf{Hueco.} Ningún signo, símbolo matemático o abreviatura existente
 lo cubre ya.
\item \textbf{Distinción.} Su forma no se confunde, escrita rápido, con otro signo
 ni con una letra.
\end{enumerate}

\subsection{Cómo se prueba}
\begin{enumerate}
\item Dibujo dos o tres formas posibles y elijo la más rápida y legible de verdad:
 cuentan el tiempo real, los cambios de dirección y si se sigue leyendo escrita
 pequeña. Menos trazos no siempre es más rápido: la corona o la estrella son un
 solo trazo con muchos giros.
\item La anoto en la leyenda personal con un código provisional (\texttt{C1},
 \texttt{C2}\ldots) y la uso dos semanas.
\item Al final, releo los apuntes de esas dos semanas sin la leyenda. Si la
 entiendo siempre y no la confundí, la registro con el siguiente código K libre.
 Si no, la descarto.
\end{enumerate}

\subsection{Reglas de estabilidad}
\begin{itemize}
\item Un código K nunca se reutiliza para otro significado.
\item Cambiar el dibujo de un signo está permitido si el significado no cambia;
 el dibujo anterior se conserva en el anexo para leer apuntes viejos.
\item Cambiar el significado exige retirar el signo y crear uno nuevo, salvo que
 el cambio solo amplíe su uso y quede escrito en el anexo de cambios.
\item Un signo que no uso en todo un ciclo es candidato a retirarse.
\end{itemize}

\subsection{Qué medir}
El cuaderno de práctica (sección 15) trae tres pruebas para repetir cada mes:
\begin{itemize}
\item \textbf{Velocidad} en tres formatos: español completo, español telegráfico
 con abreviaturas y KalaScript. Cada uno de los tres textos se escribe en los
 tres formatos, con el orden rotado, y los tiempos se comparan dentro de cada
 texto. La comparación que importa es KalaScript frente al español telegráfico,
 no frente a la prosa completa.
\item \textbf{Dictado}: tomar apuntes de alguien que lee o de una clase grabada,
 porque copiar no es lo mismo que escuchar.
\item \textbf{Relectura} a los 2 y a los 14 días: líneas no entendidas y
 relaciones mal escritas.
\end{itemize}
Si un signo aparece una y otra vez en la lista de relectura, el problema es el
signo, no mi memoria. Todavía no hay evidencia de que KalaScript sea más rápido
que el español telegráfico; estas pruebas son la forma de averiguarlo.

\subsection{Revisar las raíces con apuntes reales}
Las veintinueve raíces se eligieron sin apuntes reales, a partir de las ideas
que se repiten en todas las materias. Después de un mes de uso conviene revisarlas:
\begin{itemize}
\item Contar qué raíces uso y cuáles no. Una raíz que no aparece en un mes es
 candidata a retirarse.
\item Buscar palabras que escribo en español en casi todas las clases: son
 candidatas a raíz o a compuesto.
\item Revisar los compuestos del diccionario que nunca uso y los que invento
 una y otra vez en la leyenda personal.
\end{itemize}
Un compuesto nuevo es barato: basta con anotarlo en la leyenda. Una raíz nueva
es cara, porque hay que aprenderla; exige las mismas cuatro condiciones que un
signo nuevo.

\subsection{Ideas para más adelante}
\begin{itemize}
\item Variantes cursivas que unan un conector con la palabra siguiente sin
 levantar el lápiz.
\item Diccionarios de compuestos por materia, en una hoja aparte, sin cargar
 el diccionario común.
\item Una fuente tipográfica para escribir KalaScript en el computador.
\end{itemize}
